\documentclass[12pt]{article}
\usepackage[a4paper,margin=1in]{geometry}
\usepackage[utf8]{inputenc}
\usepackage[T1]{fontenc}
\usepackage{amsmath,amssymb}
\usepackage{textcomp}
\usepackage{parskip}
\begin{document}
\begin{center}
{\Large \textbf{2021 STEP II}}
\end{center}

\subsection*{Page 1}
\begin{flushleft}
\bigskip\textbf{Section A: Pure Mathematics}\par
1 Prove, from the identities for cos(A ± B), that cos a cos 3a \equiv 1 (cos 4a + cos 2a).\par
Find a similar identity for sin a cos 3a .\par
\textbf{(i)} Solve the equation\par
\[4 cos x cos 2x cos 3x = 1\]
\[for 0 \leq x \leq \pi.\]
\textbf{(ii)} Prove that if\par
\[tan x = tan 2x tan 3x tan 4x (\geq)\]
\[then cos 6x = or sin 4x = 0.\]
\[Hence determine the solutions of equation (\geq) with 0 \leq x \leq \pi.\]
\end{flushleft}
\subsection*{Page 2}
\begin{flushleft}
\bigskip\textbf{Question 2.} In this question, the numbers a, b and c may be complex.\par
\textbf{(i)} Let p, q and r be real numbers. Given that there are numbers a and b such that\par
\[a + b = p , a 2 + b 2 = q and a 3 + b 3 = r , (∗)\]
\[show that 3pq - p 3 = 2r .\]
\textbf{(ii)} Conversely, you are given that the real numbers p, q and r satisfy 3pq - p 3 = 2r . By\par
\[considering the equation 2x 2 - 2px + (p 2 - q) = 0, show that there exist numbers a\]
and b such that the three equations (∗) hold.\par
\textbf{(iii)} Let s, t, u and v be real numbers. Given that there are distinct numbers a, b and c\par
such that\par
2 2 2 3 3 3\par
\[a + b + c = s , a + b + c = t , a + b + c = u and abc = v ,\]
show, using part (i), that c is a root of the equation\par
\[6x 3 - 6sx 2 + 3(s 2 - t)x + 3st - s 3 - 2u = 0\]
and write down the other two roots.\par
\[Deduce that s 3 - 3st + 2u = 6v .\]
\textbf{(iv)} Find numbers a, b and c such that\par
\[a + b + c = 3 , a 2 + b 2 + c 2 = 1 , a 3 + b 3 + c 3 = -3 and abc = 2 , (∗∗)\]
and verify that your solution satisfies the four equations (∗∗).\par
\end{flushleft}
\subsection*{Page 3}
\begin{flushleft}
\bigskip\textbf{Question 3.} In this question, x, y and z are real numbers.\par
Let (x) denote the largest integer that satisfies (x) \leq x and let \{x\} denote the fractional\par
part of x, so that x = (x) + \{x\} and 0 \leq \{x\} < 1. For example, if x = 4.2 , then (x) = 4\par
\[and \{x\} = 0.2 and if x = -4.2 , then (x) = -5 and \{x\} = 0.8 .\]
\textbf{(i)} Solve the simultaneous equations\par
\textbf{(x)} + \{y\} = 4.9 ,\par
\[\{x\} + (y) = -1.4 .\]
\textbf{(ii)} Given that x, y and z satisfy the simultaneous equations\par
\[x + (y) + \{z\} = 3.9 ,\]
\[\{x\} + y + (z) = 5.3 ,\]
\textbf{(x)} + \{y\} + z = 5 ,\par
\[show that \{y\} + (z) = 3.2 and solve the equations.\]
\textbf{(iii)} Solve the simultaneous equations\par
\[x + 2(y) + \{z\} = 3.9 ,\]
\[\{x\} + 2y + (z) = 5.3 ,\]
\textbf{(x)} + 2\{y\} + z = 5 .\par
\bigskip\textbf{Question 4.} (i) Sketch the curve y = xe x , giving the coordinates of any stationary points.\par
\textbf{(ii)} The function f is defined by f(x) = xe x for x \geq a, where a is the minimum possible\par
value such that f has an inverse function. What is the value of a?\par
\[Let g be the inverse of f. Sketch the curve y = g(x).\]
\textbf{(iii)} For each of the following equations, find a real root in terms of a value of the function g,\par
or demonstrate that the equation has no real root. If the equation has two real roots,\par
determine whether the root you have found is greater than or less than the other root.\par
\[(a) e -x = 5x (b) 2x ln x + 1 = 0 (c) 3x ln x + 1 = 0 (d) x = 3 ln x\]
\textbf{(iv)} Given that the equation x x = 10 has a unique positive root, find this root in terms of\par
a value of the function g.\par
\end{flushleft}
\subsection*{Page 4}
\begin{flushleft}
\bigskip\textbf{Question 5.} (i) Use the substitution y = (x - a)u , where u is a function of x, to solve the differential\par
equation\par
dy\par
\[(x - a) = y - x ,\]
dx\par
where a is a constant.\par
\textbf{(ii)} The curve C with equation y =\leq f(x) \geqhas the property that, for all values of t except\par
\[t = 1, the tangent at the point t, f(t) passes through the point (1, t).\]
\[(a) Given that f(0) = 0, find f(x) for x < 1.\]
\[Sketch C for x < 1. You should find the co-ordinates of any stationary points and\]
consider the gradient of C as x \to 1 . You may assume that z ln |z| \to 0 as z \to 0 .\par
\[(b) Given that f(2) = 2, sketch C for x > 1, giving the co-ordinates of any stationary\]
points.\par
\end{flushleft}
\subsection*{Page 5}
\begin{flushleft}
\bigskip\textbf{Question 6.} A plane circular road is bounded by two concentric circles with centres at point O. The\par
inner circle has radius R and the outer circle has radius R + w. The points A and B lie on\par
the outer circle, as shown in the diagram, with ∠AOB = 2α, 1 \pi \leq α \leq 1 \pi and 0 < w < R.\par
3 2\par
A B\par
O\par
O\par
\textbf{(i)} Show that I cannot cycle from A to B in a straight line, while remaining on the road.\par
\textbf{(ii)} I take a path from A to B that is an arc of a circle. This circle is tangent to the inner\par
\[\leq\]
\[edge of the road, and has radius R + d (where d > w) and centre O .\]
My path is represented by the dashed arc in the above diagram.\par
\[Let ∠AO \leq B = 2θ.\]
(a) Use the cosine rule to find d in terms of w, R and cos α .\par
(b) Find also an expression for sin(α - θ) in terms of R, d and sin α .\par
w\par
You are now given that is much less than 1.\par
R\par
d\par
\textbf{(iii)} Show that and α - θ are also both much less than 1.\par
R\par
\textbf{(iv)} My friend cycles from A to B along the outer edge of the road.\par
Let my path be shorter than my friend's path by distance S. Show that\par
\[S = 2(R + d)(α - θ) + 2α(w - d).\]
Hence show that S is approximately a fraction\par
( )\par
sin α - α cos α w\par
α(1 - cos α) R\par
of the length of my friend's path.\par
\end{flushleft}
\subsection*{Page 6}
\begin{flushleft}
\bigskip\textbf{Question 7.} (i) The matrix R represents an anticlockwise rotation through angle φ (0 \textasciicircum{}\circ \leq φ < 360 \textasciicircum{}\circ )\par
in two dimensions, and the matrix R + I also represents a rotation in two dimensions.\par
\[Determine the possible values of φ and deduce that R^3 = I.\]
\textbf{(ii)} Let S be a real matrix with S3 = I, but S \sum= I.\par
\[Show that det(S) = 1.\]
Given that ( )\par
a b\par
\[S =\]
c d\par
\[show that S^2 = (a + d)S - I.\]
\[Hence prove that a + d = -1.\]
\textbf{(iii)} Let S be a real 2 \times 2 matrix.\par
\[Show that if S^3 = I and S + I represents a rotation, then S also represents a rotation.\]
What are the possible angles of the rotation represented by S?\par
\bigskip\textbf{Question 8.} (i) Show that, for n = 2, 3, 4, . . . ,\par
\[\int \sum ( )\]
d\par
\[t n (1 - t) n = nt n-2 (1 - t) n-2 (n - 1) - 2(2n - 1)t(1 - t) .\]
\[dt2\]
\textbf{(ii)} The sequence T0, T1, . . . is defined by\par
[\par
1 t n (1 - t) n\par
t\par
\[Tn = e dt .\]
n!\par
\[Show that, for n \geq 2,\]
\[Tn = Tn-2 - 2(2n - 1)Tn-1 .\]
\textbf{(iii)} Evaluate T0 and T1 and deduce that, for n \geq 0, Tn can be written in the form\par
\[Tn = an + bne ,\]
where an and bn are integers (which you should not attempt to evaluate).\par
e\par
\textbf{(iv)} Show that 0 < Tn < for n \geq 0. Given that bn is non-zero for all n, deduce that\par
n!\par
-an\par
tends to e as n tends to infinity.\par
bn\par
\end{flushleft}
\subsection*{Page 7}
\begin{flushleft}
\bigskip\textbf{Section B: Mechanics}\par
\bigskip\textbf{Question 9.} Two particles, of masses m1 and m2 where m1 > m2, are attached to the ends of a light,\par
inextensible string. A particle of mass M is fixed to a point P on the string. The string\par
passes over two small, smooth pulleys at Q and R, where QR is horizontal, so that the\par
\[particle of mass m^1 hangs vertically below Q and the particle of mass m^2 hangs vertically\]
below R. The particle of mass M hangs between the two pulleys with the section of the\par
string P Q making an acute angle of θ1 with the upward vertical and the section of the string\par
P R making an acute angle of θ2 with the upward vertical. S is the point on QR vertically\par
above P . The system is in equilibrium.\par
\textbf{(i)} Using a triangle of forces, or otherwise, show that:\par
(a)\par
]\par
\[m^2\]
\[- m^2\]
\[< M < m^1 + m^2 ;\]
(b) S divides QR in the ratio r : 1, where\par
M 2 - m 2 + m 2\par
1 2\par
\[r = .\]
\[M 2 - m^2 + m^2\]
2 1\par
\textbf{(ii)} You are now given that M 2 = m 2 + m 2 .\par
1 2\par
\textasciicircum{}\circ\par
\[Show that θ1 + θ2 = 90 and determine the ratio of QR to SP in terms of the masses\]
only.\par
\end{flushleft}
\subsection*{Page 8}
\begin{flushleft}
\bigskip\textbf{Question 10.} A train moves westwards on a straight horizontal track with constant acceleration a, where\par
\[a > 0. Axes are chosen as follows: the origin is fixed in the train; the x-axis is in the\]
direction of the track with the positive x-axis pointing to the East; and the positive y-axis\par
points vertically upwards.\par
A smooth wire is fixed in the train. It lies in the x-y plane and is bent in the shape given by\par
ky = x , where k is a positive constant. A small bead is threaded onto the wire. Initially,\par
the bead is held at the origin. It is then released.\par
\textbf{(i)} Explain why the bead cannot remain stationary relative to the train at the origin.\par
\textbf{(ii)} Show that, in the subsequent motion, the coordinates (x, y) of the bead satisfy\par
\[x˙ (x¨ - a) + y˙(y¨ + g) = 0\]
and deduce that 1 (x˙ 2 + y˙ 2 ) - ax + gy is constant during the motion.\par
\textbf{(iii)} Find an expression for the maximum vertical displacement, b, of the bead from its\par
initial position in terms of a, k and g.\par
\textbf{(iv)} Find the value of x for which the speed of the bead relative to the train is greatest and\par
give this maximum speed in terms of a, k and g.\par
\end{flushleft}
\subsection*{Page 9}
\begin{flushleft}
\bigskip\textbf{Section C: Probability and Statistics}\par
\bigskip\textbf{Question 11.} A train has n seats, where n \geq 2. For a particular journey, all n seats have been sold, and\par
each of the n passengers has been allocated a seat.\par
The passengers arrive one at a time and are labelled T1, . . . , Tn according to the order in\par
which they arrive: T1 arrives first and Tn arrives last. The seat allocated to Tr (r = 1, . . . , n)\par
is labelled Sr .\par
Passenger T1 ignores their allocation and decides to choose a seat at random (each of the n\par
seats being equally likely). However, for each r \geq 2, passenger Tr sits in Sr if it is available\par
or, if Sr is not available, chooses from the available seats at random.\par
\textbf{(i)} Let Pn be the probability that, in a train with n seats, Tn sits in Sn . Write down the\par
\[value of P^2 and find the value of P^3 .\]
\textbf{(ii)} Explain why, for k = 2, 3, . . . , n - 1,\par
\[\leq \geq\]
\[P Tn sits in Sn | T^1 sits in Sk = Pn-k+1 ,\]
\[and deduce that, for n \geq 3,\]
\{ \alpha\par
n\}-1\par
\[Pn = 1 + Pr .\]
n\par
\[r=2\]
\textbf{(iii)} Give the value of Pn in its simplest form and prove your result by induction.\par
\textbf{(iv)} Let Qn be the probability that, in a train with n seats, Tn-1 sits in Sn-1 . Determine\par
\[Qn for n \geq 2.\]
\end{flushleft}
\subsection*{Page 10}
\begin{flushleft}
\bigskip\textbf{Question 12.} (i) A game for two players, A and B, can be won by player A, with probability pA, won by\par
\[player B, with probability pB, where 0 < pA + pB < 1, or drawn. A match consists of a\]
series of games and is won by the first player to win a game. Show that the probability\par
that A wins the match is\par
pA\par
.\par
pA + pB\par
\textbf{(ii)} A second game for two players, A and B, can be won by player A, with probability p,\par
\[or won by player B, with probability q = 1 - p. A match consists of a series of games\]
and is won by the first player to have won two more games than the other. Show that\par
the match is won after an even number of games, and that the probability that A wins\par
the match is\par
p\par
.\par
\[p^2 + q^2\]
\textbf{(iii)} A third game, for only one player, consists of a series of rounds. The player starts the\par
game with one token, wins the game if they have four tokens at the end of a round and\par
loses the game if they have no tokens at the end of a round. There are two versions of\par
the game. In the cautious version, in each round where the player has any tokens, the\par
\[player wins one token with probability p and loses one token with probability q = 1-p.\]
In the bold version, in each round where the player has any tokens, the player's tokens\par
\[are doubled in number with probability p and all lost with probability q = 1 - p.\]
In each of the two versions of the game, find the probability that the player wins.\par
\[Hence show that the player is more likely to win in the cautious version if 1 > p >\]
\[and more likely to win in the bold version if 0 < p < .\]
\end{flushleft}
\end{document}
